\documentclass[10pt,article]{IEEEtran}
\usepackage{braket} 
\usepackage{cite}
\usepackage{amsmath}
\begin{document}
\section{Basic quantum information}
In quantum computing, the fundamental unit of information is the quantum bit or qubit. A qubit can exist in a superposition of states, represented as:
\begin{equation}
  \ket{\psi} = \alpha\ket{0} + \beta\ket{1},
  \label{eq:qubit_state}
\end{equation}
where \(\alpha\) and \(\beta\) are complex numbers satisfying the normalization condition \(|\alpha|^2 + |\beta|^2 = 1\). 
The states \(\ket{0}\) and \(\ket{1}\) are the computational basis states.
Another important concept is the density matrix, which provides a more general description of quantum states, especially for mixed states. The density matrix \(\rho\) for a pure state \(\ket{\psi}\) is given by:
\begin{equation}
  \rho = \ket{\psi}\bra{\psi}.
  \label{eq:density_matrix}
\end{equation}
The density matrix must satisfy the properties of being Hermitian, having a trace equal to one, and being positive semi-definite.
\section{Entanglement}
Entanglement is a unique quantum phenomenon where the states of two or more qubits become correlated such that the state of one qubit cannot be described independently of the others. For example, the Bell state is an example of a maximally entangled state:
\begin{equation}
  \ket{\Phi^+} = \frac{1}{\sqrt{2}}\left(\ket{00} + \ket{11}\right).
  \label{eq:bell_state}
\end{equation}
Entangled states are crucial for various quantum computing protocols, including quantum teleportation and superdense coding.    
\section{Quantum gates}
Quantum gates are the building blocks of quantum circuits, manipulating qubits through unitary operations. Common single-qubit gates include the Pauli-X, Pauli-Y, and Pauli-Z gates, represented by the following matrices:
\begin{equation}
    X = \begin{bmatrix}
        0 & 1 \\
        1 & 0
    \end{bmatrix}, \quad
    Y = \begin{bmatrix}
        0 & -i \\
        i & 0
    \end{bmatrix}, \quad
    Z = \begin{bmatrix}
        1 & 0 \\
        0 & -1
    \end{bmatrix}.
    \label{eq:pauli_gates}
\end{equation}
Two-qubit gates, such as the CNOT gate, are essential for creating entanglement between qubits. The CNOT gate is represented by the matrix:
\begin{equation}
    \text{CNOT} = \begin{bmatrix}
        1 & 0 & 0 & 0 \\
        0 & 1 & 0 & 0 \\
        0 & 0 & 0 & 1 \\
        0 & 0 & 1 & 0
    \end{bmatrix}.
    \label{eq:cnot_gate}
\end{equation}

Test the citation management system~\cite{Chuang2010quantumbook}.

\bibliographystyle{IEEEtran}
\bibliography{./bib/bib.bib}

\end{document}